\chapter*{Resumo}

A identificação de regiões promotoras é uma etapa importante da anotação
funcional de genomas, pois essas regiões regulam o início da transcrição gênica.
Esta tese investiga abordagens de aprendizado de máquina para predição de
promotores em sequências de DNA, com ênfase em arquiteturas baseadas em cápsulas
e em representações derivadas de propriedades físico-químicas. A pesquisa parte
dos resultados obtidos na fase de qualificação, incluindo o modelo CapsProm e o
estudo de atributos físico-químicos, e os reorganiza em um protocolo
experimental reprodutível. O trabalho considera datasets definidos por fonte derivados
de CNNProm, GraphPro, Bhandari et al. e CDBProm, normalizados para janelas de 81
pares de bases ao redor do sítio de início de transcrição. Os experimentos serão
executados com controle explícito de proveniência, especificação formal dos
experimentos e geração de resultados diretamente reutilizáveis no texto da
tese.

\noindent \textbf{Palavras-chave:} bioinformática; predição de promotores;
redes neurais de cápsulas; aprendizado de máquina; sequências de DNA.
